\documentclass[11pt,a4paper]{article}

\usepackage[utf8]{inputenc}
\usepackage[T1]{fontenc}
\usepackage{amsmath,amssymb,amsthm}
\usepackage{mathtools}
\usepackage{hyperref}
\usepackage{enumitem}
\usepackage[margin=1in]{geometry}

%% theorem environments
\newtheorem{theorem}{Theorem}[section]
\newtheorem{proposition}[theorem]{Proposition}
\newtheorem{lemma}[theorem]{Lemma}
\newtheorem{corollary}[theorem]{Corollary}
\theoremstyle{definition}
\newtheorem{definition}[theorem]{Definition}
\newtheorem{example}[theorem]{Example}
\newtheorem{computation}[theorem]{Computation}
\theoremstyle{remark}
\newtheorem{remark}[theorem]{Remark}
\newtheorem{observation}[theorem]{Observation}

%% notation
\newcommand{\NN}{\mathbb{N}}
\newcommand{\ZZ}{\mathbb{Z}}
\newcommand{\QQ}{\mathbb{Q}}
\newcommand{\factorDiff}{\mathrm{D}}
\newcommand{\fiberSet}{\mathrm{F}}
\newcommand{\BSet}{\mathrm{B}}
\newcommand{\abs}[1]{\lvert #1\rvert}

\title{Erd\H{o}s Problem E885: Formal Findings Report}
\author{Generated from the Lean formalization in\\
\texttt{FormalConjectures/Problems/Erdos/E885/}}
\date{April 2026}

\begin{document}
\maketitle

% -------------------------------------------------------------------------
\begin{abstract}
We report on the current state of a Lean~4 / Mathlib formalization
of Erd\H{o}s Problem~885, which asks whether, for every $k\ge 1$,
there exist integers $N_1<\cdots<N_k$ whose factor-difference sets
share at least $k$ common elements.
The formalization covers:
(i) exact algebraic reformulations (square-discriminant, anchored-packet,
column/cocycle/transfer laws);
(ii) verified explicit configurations for $k=2$, $k=3$, and a $3\times 4$
configuration;
(iii) structural results on pair-fiber finiteness and B-set intersections;
(iv) the guidepost rigidity theorem showing the known $3\times 4$ seed
admits exactly four common secants;
(v) lift-surface computations certifying that no new integral fifth
column arises from the divisor-parametrized search;
(vi) addendum counterexamples to several conjectured cap bounds.
All listed results compile without \texttt{sorry} in the repository.
We distinguish carefully between exact formal theorems, verified
computations, partial formalizations, and open directions.
\end{abstract}

% -------------------------------------------------------------------------
\section{Introduction and Problem Status}\label{sec:intro}

Erd\H{o}s and Rosenfeld~\cite{ErRo97} introduced the \emph{factor-difference
set} of a positive integer and asked whether arbitrarily large common
intersections can be achieved.  The question remains open:
the current literature covers $k=2$ (Erd\H{o}s--Rosenfeld~\cite{ErRo97}),
$k=3$ (Jim\'enez-Urroz~\cite{Ji99}), with further structural analysis
by Bremner~\cite{Br19}.

The Lean formalization in this repository provides:
\begin{itemize}[nosep]
\item a complete formal problem statement (\texttt{Main.lean});
\item verified witnesses for $k=2$ and $k=3$;
\item a verified $3\times 4$ configuration (three integers sharing four
      common factor differences);
\item a library of exact algebraic identities and structural lemmas
      (\texttt{Findings/} directory);
\item finite rigidity and lift-surface computations;
\item addendum counterexamples to several cap conjectures.
\end{itemize}
The problem itself (\texttt{erdos\_885}) is stated with
\texttt{answer(sorry)}, reflecting its open status.

% -------------------------------------------------------------------------
\section{Problem Statement and Notation}\label{sec:statement}

\begin{definition}[Factor-difference set; \texttt{Defs.lean}]\label{def:factordiff}
For $n\in\NN$, $n\ge 1$, the \emph{factor-difference set} of~$n$ is
\[
  \factorDiff(n) \;=\; \bigl\{\,\abs{a-b} : a,b\in\NN,\; a,b>0,\;
  a\cdot b = n\bigr\}.
\]
\end{definition}

\begin{definition}[Fiber set; \texttt{Defs.lean}]\label{def:fiber}
For a set $S\subseteq\NN$, the \emph{fiber set} is
$\fiberSet(S)=\{n\ge 1 : S\subseteq\factorDiff(n)\}$.
\end{definition}

\begin{theorem}[Erd\H{o}s Problem 885; \texttt{Main.lean}]\label{thm:e885}
Is it true that for every $k\ge 1$ there exists a finite set
$\{N_1,\dots,N_k\}\subset\NN$ with $N_i\ge 1$ for all~$i$ and
\[
  \Bigl\lvert\bigcap_{i=1}^{k}\factorDiff(N_i)\Bigr\rvert\;\ge\; k\;?
\]
The formalization leaves the Boolean answer as \texttt{sorry} (open problem).
\end{theorem}

% -------------------------------------------------------------------------
\section{Literature Baseline: Verified Witnesses}\label{sec:witnesses}

\subsection{The case $k=2$}

\begin{theorem}[\texttt{variants.k\_eq\_2}; \texttt{Main.lean}]
$\factorDiff(12)\cap\factorDiff(42)\supseteq\{1,11\}$, witnessed by
$12 = 3\times 4$ (difference~$1$), $12 = 1\times 12$ (difference~$11$),
$42 = 6\times 7$ (difference~$1$), $42 = 3\times 14$ (difference~$11$).
Hence $\lvert\factorDiff(12)\cap\factorDiff(42)\rvert\ge 2$.
\end{theorem}

\subsection{The case $k=3$}

\begin{theorem}[\texttt{variants.k\_eq\_3}; \texttt{Main.lean}]
$\factorDiff(112)\cap\factorDiff(952)\cap\factorDiff(3240)\supseteq\{6,54,111\}$.
\end{theorem}

The individual memberships are verified by explicit factorizations in
\texttt{Computations.lean}, e.g.\ $112 = 8\times 14$ gives
$14-8=6\in\factorDiff(112)$.

\subsection{The $3\times 4$ configuration}

\begin{theorem}[\texttt{variants.three\_by\_four\_configuration}; \texttt{Main.lean}]\label{thm:3x4}
The triple $(13608,\,29848,\,65968)$ shares at least four common factor
differences $\{18,282,477,1122\}$.  Thus $k=3$ integers can achieve
$\lvert\bigcap\factorDiff(N_i)\rvert\ge 4$.
This is the strongest explicit witness in the repository; the literature
$k=4$ theorem is cited but not yet formalized.
\end{theorem}

% -------------------------------------------------------------------------
\section{Square-Discriminant Reformulation}\label{sec:sq-disc}

\begin{proposition}[Proposition~1.2; \texttt{SquareDiscriminant.lean}]\label{prop:sq-disc}
For $d,n\ge 1$,
\[
  d\in\factorDiff(n) \;\iff\;
  \exists\, m\in\NN,\; m>d,\; m^2 = d^2 + 4n.
\]
Equivalently, $d^2+4n$ is a perfect square.
\end{proposition}

Both directions are formalized separately as
\texttt{factorDiff\_implies\_square} (forward) and
\texttt{square\_implies\_factorDiff} (backward), in addition to the
biconditional \texttt{square\_discriminant\_iff}.

% -------------------------------------------------------------------------
\section{Pair-Fiber Structure and Finiteness}\label{sec:pairfiber}

\begin{theorem}[Finiteness of two-element fibers; \texttt{PairFiber.lean}]\label{thm:pairfiber-fin}
For distinct $d,e>0$, the fiber set $\fiberSet(\{d,e\})$ is finite.
\end{theorem}

This is a weakened consequence of the full pair-fiber parametrization
from the notes.  The proof uses the square-discriminant witnesses
$(x,y)$ with $x^2=d^2+4n$ and $y^2=e^2+4n$; then $y^2-x^2=e^2-d^2$
is fixed, so the factorizations of $e^2-d^2$ bound both $x$ and $y$,
hence~$n$.

\begin{theorem}[Pair-fiber parametrization; \texttt{PairFiberParametrization.lean}]
If $d<e$ are both factor differences of~$n$, the square-discriminant
witnesses satisfy $(y-x)(y+x)=e^2-d^2$, and there is a uniform bound
$n\le (e^2-d^2)^2$.
\end{theorem}

\begin{remark}
The full same-parity counting formula from the notes is \emph{not}
formalized; only finiteness and the parametrization identity are
established.
\end{remark}

% -------------------------------------------------------------------------
\section{Anchored Packets, Column Law, and Transfer Laws}\label{sec:packet}

\subsection{Anchored packet identity}

\begin{proposition}[Proposition~4.1; \texttt{ColumnLaw.lean}]
$c = r(2x+r)$ if and only if $(x+r)^2 = x^2+c$.
\end{proposition}

This is the algebraic core of the anchored-packet formulation:
a set of factor differences $\{d_1,\dots,d_r\}$ has a common support~$N$
if and only if there exists $x>0$ with $(x+r_i)^2 = x^2+c_i$ for each row,
where $c_i = d_i^2-d_1^2$.

\subsection{Column law and cocycle law}

\begin{theorem}[Column law; \texttt{ColumnLaw.lean}]\label{thm:col-law}
For a fixed column anchor~$x$ and rows $c_i = r_i(2x+r_i)$,
$c_j(2x+r_j)$,
\[
  c_j\, r_i \;-\; c_i\, r_j \;=\; r_i\, r_j\,(r_j - r_i).
\]
\end{theorem}

\begin{theorem}[Cocycle law; \texttt{ColumnLaw.lean}]
For three rows $i<j<k$,
\[
  \Delta_{jk}\,\delta_{ij} - \Delta_{ij}\,\delta_{jk}
  = \delta_{ij}\,\delta_{jk}\,\delta_{ik},
\]
where $\delta_{ij}=r_j-r_i$ and $\Delta_{ij}=c_j-c_i$.
\end{theorem}

\subsection{Three-row transfer law}

\begin{theorem}[Theorem~4.8; \texttt{ColumnLaw.lean}]\label{thm:3row-transfer}
Setting $u=r_3-r_2$, $v=r_4-r_3$, $B=c_3-c_2$, $A=c_4-c_3$,
\[
  A\,u - B\,v = u\,v\,(u+v).
\]
\end{theorem}

\subsection{Four-row transfer laws}

\begin{theorem}[\texttt{FourRowTransfer.lean}]
The three-row cubic transfer law holds on both consecutive triples of a
four-row anchored packet.  In addition, the quartic discriminant formulas
\[
  (2uv+v^2-A)^2 = v^4 - 2(A+2B)v^2 + A^2
\]
and
\[
  (2wv+v^2+A)^2 = v^4 + 2(A+2C)v^2 + A^2
\]
are established, together with the Riemann--Hurwitz arithmetic
$2\cdot 5 - 2 = 4(-2)+16$.
\end{theorem}

\subsection{Common-parameter equivalence}

\begin{theorem}[\texttt{ObstructionDescent.lean}]\label{thm:common-param}
Over~$\QQ$, the three-row transfer system
$Au - Bv = uv(u+v)$ and $Cv - Aw = vw(v+w)$ is equivalent to the
existence of a single affine parameter~$\lambda$ such that
$B = u(u+\lambda)$, $B+A = (u+v)(u+v+\lambda)$,
$B+A+C = (u+v+w)(u+v+w+\lambda)$.
\end{theorem}

\subsection{Verified $3\times 4$ packet}

The packet with row offsets $(c_2,c_3,c_4)=(79200,227205,1258560)$
and column anchors $x\in\{18,234,346,514\}$ is fully verified:
all twelve identities $c_i = r_{i,t}(2x_t + r_{i,t})$ are checked
in \texttt{ColumnLaw.lean}.

\subsection{Cross-ratio test}

\begin{theorem}[Theorem~4.5; \texttt{ColumnLaw.lean}]
The cross-ratios of the three rows of the $3\times 4$ packet are
pairwise distinct:
$9840/6144 \ne 23760/14688 \ne 52992/32256$,
hence the packet is not of M\"obius type.
\end{theorem}

% -------------------------------------------------------------------------
\section{B-Set and Spectrum Viewpoint}\label{sec:bset}

\begin{definition}[B-set; \texttt{BSet.lean}]
$\BSet(\delta) = \{2z : z>0,\; \exists\, r,\; r^2 = z^2+\delta\}$.
\end{definition}

\begin{theorem}[B-set membership; \texttt{BSet.lean}]
$2z\in\BSet(\delta)$ iff $z^2+\delta$ is a perfect square, iff
$\delta = x(x+2z)$ for some $x>0$.
\end{theorem}

\begin{theorem}[Pairwise B-set finiteness; \texttt{BSet.lean}]
For $0<a<b$, $\BSet(a)\cap\BSet(b)$ is finite.
\end{theorem}

The proof reduces to the factorizations of $b-a$, analogous to the
pair-fiber argument.  The full iff-parametrization by divisors of $b-a$ is
not encoded.

\subsection{Spectral identities}

The file \texttt{Spectrum.lean} records the column-spectrum
factorization $r^2-d^2 = (r-d)(r+d)$ and verifies:
\begin{itemize}[nosep]
\item the three $\delta_t$ values $(54432, 119392, 263872)$ used in
      the spectral-cap example;
\item four elements of $\BSet(54432)$:
      $18^2+54432=234^2$, $282^2+54432=366^2$,
      $477^2+54432=531^2$, $1122^2+54432=1146^2$;
\item the $7$-interface near-miss at $(48960,196560)$;
\item quartic genus formula $\deg/2-1 = 1$ and degree-$8$
      hyperelliptic genus~$3$.
\end{itemize}
It does \emph{not} formalize the full anchored-spectrum theorem or
spectral-cap obstruction.

% -------------------------------------------------------------------------
\section{Row-Complex and Terminal-Spectrum Viewpoints}\label{sec:rowcomplex}

\subsection{Row-complex identities (\texttt{RowComplex.lean})}

\begin{theorem}[Row-complex identity]
If $r_1^2-d_1^2=4N$ and $r_2^2-d_2^2=4N$, then
$r_2^2-r_1^2 = d_2^2-d_1^2$.
\end{theorem}

Additional verified chain computations cover Examples~6.3--6.6 from the
notes, including explicit supporting factorizations for the pairs
$\{1,11\}$, $\{6,54,111\}$, $\{2,37,58\}$, and the four-difference
packet.

\begin{theorem}[Terminal stitch identity; \texttt{RowComplex.lean}]
If $uv=\Lambda$ and $v-u=s$, then $s^2+4\Lambda = (u+v)^2$.
\end{theorem}

\subsection{Terminal spectrum (\texttt{TerminalSpectrum.lean})}

\begin{definition}
The \emph{terminal spectrum} of a column $C$ is
$\{\Lambda\in\ZZ : \forall a,\; C(a)^2+4\Lambda \text{ is a perfect square}\}$.
\end{definition}

The file establishes:
\begin{itemize}[nosep]
\item Prefix monotonicity: deleting the last row can only enlarge the
      terminal spectrum.
\item Positive spectral characterization: $\Lambda>0$ lies in the
      spectrum of a positive column iff $\Lambda=x(x+c)$ for some $x>0$.
\item Backward reconstruction: for $-\Delta$ in the spectrum,
      $U(a)=\sqrt{C(a)^2-4\Delta}$ is the unique nonneg coherent column,
      it preserves square-step profiles, and it translates the spectrum by
      exactly~$\Delta$.
\item An $8$-row, $3$-column chain example with shifts
      $\Delta\in\{786240,\,982080\}$ verified entry-by-entry.
\end{itemize}

% -------------------------------------------------------------------------
\section{Square-Translate Geometry}\label{sec:sq-translate}

\begin{theorem}[Pairwise factorization; \texttt{SquareTranslate.lean}]
Over~$\QQ$, $(H/u_i-u_i)^2-(H/u_j-u_j)^2 = (u_j^2-u_i^2)(H^2/(u_i^2u_j^2)-1)$.
\end{theorem}

\begin{theorem}[Three-shift Weierstrass model; \texttt{SquareTranslate.lean}]
Fix shifts $(0,p,q)$.  The common-square curve
$v^2=u^2+p$, $w^2=u^2+q$ is birationally equivalent to the elliptic curve
$E_{p,q}: Y^2 = X(X+4p)(X+4q)$
via $X=4(u+v)(u+w)$, $Y=8(u+v)(u+w)(v+w)$.
The inverse map is verified as $16pq-X^2 = -4Yu$.
\end{theorem}

\begin{theorem}[Genus formulas; \texttt{SquareTranslate.lean}]
The static-support genus evaluates to $0$ ($r=2$), $1$ ($r=3$), $5$ ($r=4$),
confirming that three shifts give an elliptic curve and four shifts give
genus~$5$.
\end{theorem}

\begin{example}[Normalized seeds; \texttt{SquareTranslate.lean}]
The triple $(112,952,3240)$ normalizes to shifts $e=(0,2880,12285)$
with $u$-values $(6,22,62)$; the quadruple normalizes to
$e=(0,79200,227205,1258560)$ with $u=(18,234,346,514)$.
Both sets of square identities are formally verified.
\end{example}

\subsection{Non-split seed analysis}

The seed $(79200,227205,1258560)$ is studied further:
\begin{itemize}[nosep]
\item Product factorization $(r-p)(r-q)=1179360\times 1031355$.
\item Four rational points on the seed curve at $u\in\{18,234,346,514\}$.
\item Three points are collinear on $Y=1582X+17319744$.
\item All four lie on the elliptic curve
      $Y^2=X(X+316800)(X+908820)$.
\item The cross-ratio $\lambda=227205/79200=459/160$ is not a perfect
      square in~$\QQ$ (formalized with a denominator argument).
\item Deformation base verification:
      $\Delta^2=(m-1532)(m-524)(m+596)(m+1460)$ at $m=1582$.
\end{itemize}

% -------------------------------------------------------------------------
\section{Shift-Transpose}\label{sec:shift-transpose}

\begin{theorem}[\texttt{ShiftTranspose.lean}]
The transpose identity: if $x_0^2=y_0^2+a$ and $z^2=y^2+a$, then
$x_0^2+(y^2-y_0^2)=z^2$.  This immediately yields a square witness
for the translated shift and underlies the passage from a
$5$-point/$3$-shift configuration to a $4$-point/$4$-shift one.
\end{theorem}

% -------------------------------------------------------------------------
\section{Divisor Templates and Projective Obstructions}\label{sec:divisor}

\begin{theorem}[Factor-pair encoding; \texttt{DivisorTemplate.lean}]
If $a=r(r+n)$, $b=s(s+n)$, $c=t(t+n)$, and $M=2s+n$, then
$b-a = x(M-x)$, $c-b = y(M+y)$, and $M^2-4b=n^2$,
where $x=s-r$, $y=t-s$.
\end{theorem}

\begin{theorem}[Divisor-template identity; \texttt{DivisorTemplate.lean}]
$y^2-A = a+u^2$ and $y^2+B = c+u^2$ where $y^2=b+u^2$, $A=b-a$, $B=c-b$.
\end{theorem}

\begin{theorem}[M\"obius obstruction; \texttt{DivisorTemplate.lean}]
Both the diagonal and antidiagonal projective-rank-$2$ obstructions
force $\kappa=\kappa'$: a nontrivial functional identity
$\kappa z - \lambda^2 z^3 = \kappa'\lambda z - \lambda z^3$ for all
$z\ne 0$ implies $\kappa=\kappa'$.
\end{theorem}

% -------------------------------------------------------------------------
\section{Obstruction Descent}\label{sec:descent}

\begin{theorem}[Factor-difference transfer; \texttt{ObstructionDescent.lean}]
If $x^2=d^2+4n$ and $t^2=d^2+4M$, then $t\in\factorDiff(n)$ iff
$x\in\factorDiff(M)$.  The full incidence matrix of common
factor-difference memberships is transported exactly to the lower-level
matrix.
\end{theorem}

\begin{theorem}[Anchored three-shift reduction; \texttt{ObstructionDescent.lean}]
The three-shift system $z_1^2=u^2+p$, $z_2^2=u^2+q$, $z_3^2=u^2+r$
with all entries positive is equivalent to the rigid three-square
system $M^2-4p=v_1^2$, $M^2-4q=v_2^2$, $M^2-4r=v_3^2$
where $M$ and the $v_i$ have prescribed parities.
\end{theorem}

% -------------------------------------------------------------------------
\section{Guidepost Rigidity}\label{sec:guidepost}

\begin{theorem}[\texttt{GuidepostRigidity.lean}]\label{thm:guidepost}
The seed triple $(79200,\,227205,\,1258560)$ has exactly four
positive common secants: $\{36,\,468,\,692,\,1028\}$.  In particular,
the known $3\times 4$ seed is \emph{rigid}: no fifth common secant exists
for this row triple.
\end{theorem}

The proof proceeds by:
\begin{enumerate}[nosep]
\item Computing the sum image
      $\Sigma(148005)$ (24~elements, via \texttt{native\_decide}).
\item Computing the difference image
      $\Delta(1031355)$ (24~elements).
\item Intersecting to obtain the four middle values
      $\{954,\,1062,\,1178,\,1402\}$.
\item Applying the square test $d^2 = m^2 - 4b$ to recover the four
      secants $\{36,468,692,1028\}$.
\item Proving the iff characterization:
      $d\in\factorDiff(A)\cap\factorDiff(B)\cap\factorDiff(C)$
      with $d>0$ iff $d\in\{36,468,692,1028\}$.
\end{enumerate}

% -------------------------------------------------------------------------
\section{Lift-Surface Computations}\label{sec:lift}

\begin{definition}[\texttt{LiftSurface.lean}]
The \emph{moving lift parameter} $m(u,v)$ for the packet
$(p,q,r)=(79200,227205,1258560)$ is the rational function
\[
  m(u,v) = \frac{(u^2-324)v^2-240552\,uv-324\,u^2+14569656624}
                {18(uv-120276)}.
\]
The base quartic is
$(m-1532)(m-524)(m+596)(m+1460)$.
\end{definition}

\begin{theorem}[Centered factorizations; \texttt{LiftSurface.lean}]
Each linear factor of the quartic, multiplied by $18(uv-120276)$,
factors as a product of two bilinear forms in $(u,v)$.  The product of
all four factorizations yields
\[
  \Delta(m(u,v))\cdot 18^4(uv-120276)^4 = \Phi(u,v),
\]
where $\Phi$ is an explicit degree-$8$ polynomial.
\end{theorem}

\begin{computation}[Rational lift point; \texttt{LiftSurface.lean}]
$(u,v)=(3406,18)$ gives $m=6548/39$ (non-integral),
$\Delta(m)=(1182146560/1521)^2$.
\end{computation}

\begin{theorem}[Finite search; \texttt{LiftSurface.lean}]
The divisor parametrization of $r=1258560$ yields $60$ positive
absolute lift abscissas.  The square test on admissible signed pairs
produces exactly six solutions:
\[
  \{(-3406,-18),\;(-514,234),\;(-234,514),\;(234,-514),\;(514,-234),\;(3406,18)\}.
\]
Every surviving pair either has non-integral~$m$ or reuses a built-in
fifth value $\abs{v}\in\{18,234,514\}$.  Hence no genuine integral fifth
column arises from this search.
\end{theorem}

% -------------------------------------------------------------------------
\section{Addendum Computations and Counterexamples}\label{sec:addendum}

The file \texttt{AddendumComputations.lean} records several further
explicit configurations that were discovered after the initial
formalization and serve as counterexamples to various conjectured cap
bounds.

\begin{computation}[$3$-row / $5$-column packet]
An explicit $3\times 5$ packet exists, refuting the universal bound
$\lvert Y(a,b,c)\rvert\le 4$.  The $15$ square identities
(e.g.\ $330^2+756000=930^2$) are verified by \texttt{native\_decide}.
\end{computation}

\begin{computation}[Two-shift cap counterexample]
A two-shift configuration with six solutions (twelve square identities)
is verified, refuting certain two-shift cap conjectures.
\end{computation}

\begin{computation}[$3$-point obstruction counterexample]
The set $\{18,66,186\}$ has at least four positive shifts,
refuting the $3$-point obstruction route.
\end{computation}

\begin{computation}[Primitive $4$-secant triples]
Two primitive maximal $4$-secant triples are verified:
\begin{enumerate}[nosep]
\item Supports $(46592,\,363545,\,1498112)$ with secants
      $(16,344,776,1424)$.
\item Supports $(24640,\,405405,\,1723680)$ with secants
      $(108,516,1212,1524)$.
\end{enumerate}
\end{computation}

% -------------------------------------------------------------------------
\section{Summary of Formalization Status}\label{sec:status}

\subsection{Exact formal theorems (no \texttt{sorry})}

\begin{enumerate}[nosep]
\item Square-discriminant reformulation (Proposition~\ref{prop:sq-disc}).
\item Pair-fiber finiteness and parametrization (Section~\ref{sec:pairfiber}).
\item B-set membership characterization and pairwise finiteness.
\item Column law, cocycle law, three- and four-row transfer laws.
\item Common-parameter equivalence (Theorem~\ref{thm:common-param}).
\item Factor-difference transfer and matrix transport.
\item Anchored three-shift reduction.
\item Divisor-template and M\"obius obstruction identities.
\item Shift-transpose identity and forward witness.
\item Three-shift Weierstrass model and inverse map.
\item Genus formulas for $r=2,3,4$.
\item Cross-ratio non-M\"obius test.
\item Terminal-spectrum structural theory (backward reconstruction,
      spectrum translation, prefix monotonicity, constant-shift
      propagation).
\item Guidepost rigidity: exact $4$-element common secant set
      (Theorem~\ref{thm:guidepost}).
\item Lift-surface: centered factorizations, quartic identity, finite
      search yielding six degenerate solutions.
\item All verified witnesses ($k=2$, $k=3$, $3\times 4$).
\end{enumerate}

\subsection{Verified explicit computations}

All arithmetic checks listed in Section~\ref{sec:witnesses},
\ref{sec:guidepost}, \ref{sec:lift}, and~\ref{sec:addendum} compile
without \texttt{sorry}.  Many use \texttt{native\_decide} or
\texttt{norm\_num}.

\subsection{Partial formalizations}

\begin{enumerate}[nosep]
\item \textbf{Pair-fiber counting formula}: only finiteness is proved;
      the exact same-parity divisor count is not formalized.
\item \textbf{B-set intersection parametrization}: only finiteness;
      the full iff-by-divisors is not encoded.
\item \textbf{Anchored spectrum}: only the algebraic identity and
      selected numerical checks; the full theorem, boundary embedding,
      and spectral-cap obstruction are not formalized.
\item \textbf{Row-complex}: chain computations and the key identity are
      present, but the full $H(T)/S(T)$ recursion and row-complex
      bijection are not.
\item \textbf{Square-translate geometry}: birational maps and genus
      formulas are present, but the full algebraic-geometry statements
      (e.g.\ Mordell--Weil rank, rational-point structure) are not.
\end{enumerate}

\subsection{Open directions}

\begin{enumerate}[nosep]
\item The problem itself remains open (\texttt{erdos\_885} has
      \texttt{answer(sorry)}).
\item The literature $k=4$ result is not yet formalized.
\item No $k\ge 5$ witness is known.
\item The spectral-cap and global-cap obstructions are not formalized,
      and the addendum counterexamples show some conjectured cap bounds
      are false.
\item The Mordell--Weil analysis of the three-shift elliptic curve and
      the genus-$5$ four-shift curve are not formalized.
\item The lift-surface analysis is complete for the specific seed
      $(79200,227205,1258560)$ but has not been generalized.
\end{enumerate}

% -------------------------------------------------------------------------
\section*{References}

\begin{thebibliography}{9}
\bibitem{ErRo97}
P.~Erd\H{o}s and M.~Rosenfeld,
\emph{The factor-difference set of integers},
Acta Arith.\ \textbf{79} (1997), 353--359.

\bibitem{Ji99}
J.~Jim\'enez-Urroz,
\emph{A note on a conjecture of Erd\H{o}s and Rosenfeld},
J.~Number Theory \textbf{77} (1999), 64--69.

\bibitem{Br19}
A.~Bremner,
\emph{On a problem of Erd\H{o}s related to common factor differences},
Integers \textbf{19} (2019), \#A32.

\bibitem{erdosproblems}
\url{https://www.erdosproblems.com/885}
\end{thebibliography}

\end{document}
